% =====================================================================
% Frente 4 --- Agrupamento e comunidades (Arthur)
% =====================================================================
\section{Frente 4 --- Agrupamento e comunidades}

\noindent\textbf{Responsável:} Arthur.
\textbf{No relatório:} comunidades; limitações e ética.
\textbf{Na Etapa 1:} diagrama de atividades do Louvain.

\subsection{O que foi implementado}

{\small\begin{tabularx}{\linewidth}{@{}L{3.9cm} L{4.7cm} Y L{1.9cm}@{}}
\toprule
Função & Arquivo & Complexidade & Situação \\
\midrule
\codigo{triangulos} & \codigo{metricas/agrupamento.py} & $O(m^{1,5})$~\cite{schank2005} & \afazer \\
\codigo{agrupamento\_local}, \codigo{agrupamento\_medio} & \codigo{metricas/agrupamento.py} & \falta & \afazer \\
\codigo{transitividade} & \codigo{metricas/agrupamento.py} & \falta & \afazer \\
\codigo{modularidade} & \codigo{comunidades/} & $O(m)$~\cite{newman2004} & \afazer \\
\codigo{louvain} & \codigo{comunidades/} & $\approx O(m)$ por passada~\cite{blondel2008} & \afazer \\
\codigo{propagacao\_de\_rotulos} & \codigo{comunidades/} & $O(m)$ por iteração~\cite{raghavan2007} & \afazer \\
\codigo{nmi}, \codigo{ari} & \codigo{comunidades/} & \pendente{} ~\cite{danon2005,hubert1985} & \afazer \\
\bottomrule
\end{tabularx}}

\subsection{Decisões e justificativas}
\begin{itemize}[leftmargin=*]
  \item \pendente{ex.: fórmula do ganho de modularidade usada no Louvain}
\end{itemize}

\subsection{Corretude: testes e resultados}

{\small\begin{tabularx}{\linewidth}{@{}Y L{2.5cm} L{2.2cm} L{3.8cm}@{}}
\toprule
Teste & Base & Obtido & Referência \\
\midrule
Clustering médio, transitividade e triângulos & Karate & \falta & 0,5706; 0,2557; 45 \\
Modularidade do Louvain & Karate & \falta & entre 0,39 e 0,42 (máximo 0,4198) \\
Modularidade da divisão real (gabarito) & Karate & \falta & 0,3582 \\
Clustering médio, transitividade e triângulos & Twitch PT & \falta & 0,3199; 0,131; 173.510 \\
Louvain & Twitch PT & \falta & 0,286 a 0,296 ($\approx$ 7 comunidades) \\
Clustering médio, transitividade e triângulos & ego-Facebook & \falta & 0,6055; 0,5192; 1.612.010 \\
NMI e ARI de partições iguais & sintético & \falta & 1 e 1 \\
\bottomrule
\end{tabularx}}

\subsection{Figuras e tabelas para o relatório}
\begin{itemize}[leftmargin=*]
  \item \pendente{ex.: diagrama de atividades do Louvain; Karate colorido por comunidade}
\end{itemize}

\subsection{Limitações, ética e pendências}
\begin{itemize}[leftmargin=*]
  \item Limite de resolução da modularidade~\cite{fortunato2007}. \pendente{discutir}
  \item LGPD e anonimização dos dados. \pendente{discutir}
\end{itemize}

\subsection{Registro}
\begin{description}[leftmargin=2.6cm,style=nextline]
  \registro{AAAA-MM-DD}{\pendente{primeira entrada}}
\end{description}
